
\documentclass{article}
\usepackage{colortbl}
\usepackage{makecell}
\usepackage{multirow}
\usepackage{supertabular}

\begin{document}

\newcounter{utterance}

\centering \large Interaction Transcript for game `textmapworld', experiment `medium\_cycle', episode 32 with reka{-}flash{-}3.1{-}t0.0.
\vspace{24pt}

{ \footnotesize  \setcounter{utterance}{1}
\setlength{\tabcolsep}{0pt}
\begin{supertabular}{c@{$\;$}|p{.15\linewidth}@{}p{.15\linewidth}p{.15\linewidth}p{.15\linewidth}p{.15\linewidth}p{.15\linewidth}}
    \# & $\;$A & \multicolumn{4}{c}{Game Master} & $\;\:$B\\
    \hline

    \theutterance \stepcounter{utterance}  
    & & \multicolumn{4}{p{0.6\linewidth}}{
        \cellcolor[rgb]{0.9,0.9,0.9}{
            \makecell[{{p{\linewidth}}}]{
                \texttt{\tiny{[P1$\langle$GM]}}
                \texttt{Please help me with the following task. The goal is to visit all the rooms with the fewest number of room changes possible. In each room, you need to decide the direction to go in. Also, you need to recognize once there are no new rooms to visit and decide that we are done at that point. Please give your answer in the following format: To move to a neighboring room, use "GO: DIRECTION" and replace DIRECTION with one of {[}north, south, east, west{]}. To stop the exploration, answer with "DONE" instead. Omit any other text.} \\
\texttt{Here is an example:} \\
\texttt{You are in the Kitchen. Currently available directions: south, west. What is your next instruction?} \\
\texttt{GO: west} \\
\texttt{You have made a step and entered a Lobby. Currently available directions: east, north. What is your next instruction?} \\
\texttt{GO: north} \\
\texttt{...} \\
\texttt{You have made a step and entered a Bedroom. Currently available directions: south. What is your next instruction?} \\
\texttt{DONE} \\
\texttt{Let us start. You are in the Utility room. Currently available directions: east. What is your next instruction?} \\
            }
        }
    }
    & \\ \\

    \theutterance \stepcounter{utterance}  
    & \multicolumn{4}{p{0.6\linewidth}}{
        \cellcolor[rgb]{0.9,0.9,0.9}{
            \makecell[{{p{\linewidth}}}]{
                \texttt{\tiny{[P1$\rangle$GM]}}
                \texttt{<reasoning>} \\
\texttt{Let's tackle this problem. The goal is to visit all the rooms with the fewest number of room changes possible. Each time I'm in a room, I have to choose a direction to go, and I need to make sure that I don't revisit rooms unnecessarily. Also, once there are no new rooms left to visit, I should stop.} \\
\\ 
\texttt{First, I need to understand the structure of the rooms. Since the example starts in the Kitchen and moves west to the Lobby, then north to another room, etc., it seems like the rooms are connected in a maze{-}like structure. The challenge is to find the shortest path that covers all rooms. This reminds me of the classic "traversal" problem, similar to the Hamiltonian path problem, but here we want the minimal number of steps (room changes) to visit all rooms.} \\
\\ 
\texttt{Wait, but the problem says "fewest number of room changes possible". So each move to a new room counts as a change. So the goal is to visit all rooms with the least number of moves. That's equivalent to finding the shortest possible path that visits every room at least once. This is known as the Traveling Salesman Problem (TSP), which is NP{-}hard. However, since the problem is presented in a maze{-}like scenario, maybe the structure is such that a depth{-}first search (DFS) or breadth{-}first search (BFS) approach can be used to find the minimal path.} \\
\\ 
\texttt{But how do I model the rooms and their connections? Since the user is providing the available directions each time, I need to keep track of the rooms I've visited and the paths taken. The key is to explore in a way that covers all rooms with the least steps. } \\
\\ 
\texttt{Let me think step by step. The starting point is the Utility room with only east direction available. So the first move is GO: east. Let's assume that takes me to another room, say Room A. Now, in Room A, whatever directions are available, I need to choose the one that leads to the least explored path. But how to track which rooms have been visited?} \\
\\ 
\texttt{Wait, the problem states that once there are no new rooms to visit, we stop. So the process is to keep moving until all reachable rooms are visited. But the challenge is to do this with the fewest room changes. So it's not just about visiting all rooms, but doing so in the shortest possible path. } \\
\\ 
\texttt{This sounds like the problem} \\
            }
        }
    }
    & & \\ \\

    \theutterance \stepcounter{utterance}  
    & & & \multicolumn{2}{p{0.3\linewidth}}{
        \cellcolor[rgb]{0.9,0.9,0.9}{
            \makecell[{{p{\linewidth}}}]{
                \texttt{\tiny{[GM$|$GM]}}
                \texttt{abort game} \\
            }
        }
    }
    & & \\ \\

\end{supertabular}
}

\end{document}
